\documentclass{article}
\usepackage[utf8]{inputenc}
\usepackage{graphicx}
\usepackage{listings}
\usepackage{geometry}
\geometry{a4paper, margin=1in}

\title{Manual de Usuario del Compilador HULK}
\author{David S\'anchez}
\date{\today}

\begin{document}

\maketitle

\tableofcontents

\newpage

\section{Introducción}

Este documento es una guía de usuario para el compilador del lenguaje de programación HULK. El objetivo es proporcionar instrucciones claras sobre cómo escribir código en HULK, cómo compilarlo y cómo ejecutarlo.

\section{Instalación y Compilación}

Para utilizar el compilador, necesitarás tener Flex y Bison instalados, así como un compilador de C++ (como g++). El proyecto se compila utilizando un `makefile`.

\subsection{Comandos del Makefile}

El `makefile` proporcionado automatiza las tareas de compilación y limpieza. Los principales comandos son:

\begin{itemize}
    \item \textbf{Limpiar el proyecto (`make clean`)}: Elimina todos los ficheros objeto (`.o`) y el ejecutable final. Es útil para realizar una compilación desde cero.
    \begin{verbatim}
make clean
    \end{verbatim}

    \item \textbf{Compilar el proyecto (`make compile`)}: Compila todo el código fuente y genera el ejecutable `hulk` en el directorio `hulk/`.
    \begin{verbatim}
make compile
    \end{verbatim}

    \item \textbf{Ejecutar el compilador (`make execute`)}: Este comando ejecuta el compilador. Si el ejecutable no existe o el código fuente ha cambiado, primero compilará el proyecto y luego lo ejecutará.
    \begin{verbatim}
make execute
    \end{verbatim}
\end{itemize}

\subsection{Compilar y Ejecutar un Programa HULK}

Una vez que el compilador `hulk` ha sido generado, puedes usarlo para compilar y ejecutar tus propios ficheros de código HULK (con extensión `.hulk`).

Para compilar y ejecutar un fichero, usa el siguiente comando:

\begin{verbatim}
./hulk/hulk < ruta/a/tu/fichero.hulk
\end{verbatim}

El compilador procesará el fichero y, si no hay errores, ejecutará el código.

\section{Sintaxis del Lenguaje HULK}

A continuación se presenta un resumen de la sintaxis del lenguaje HULK.

\subsection{Estructura General}

Un programa en HULK consiste en una secuencia de declaraciones y expresiones. Las expresiones deben terminar con un punto y coma (`;`).

\begin{verbatim}
// Ejemplo de un programa simple
let a = 42 in print(a);
\end{verbatim}

\subsection{Tipos de Datos}

HULK soporta los siguientes tipos de datos básicos:
\begin{itemize}
    \item \textbf{Number}: Números, tanto enteros como decimales.
    \item \textbf{String}: Cadenas de texto.
    \item \textbf{Boolean}: Valores `true` o `false`.
\end{itemize}

También soporta la definición de nuevos tipos (clases).

\subsection{Gramática}

A continuación se muestra la gramática del lenguaje en formato BNF:

\begin{verbatim}
input -> line
    | lines_block

lines_block -> { lines }

lines -> epsilon
    | non_empty_lines

non_empty_lines -> line
    | non_empty_lines line

line -> expr ;
    | func_asign ;
    | type_node_decl

expr -> arit_op
    | bool_expr
    | while_expr
    | string
    | id_expr
    | func_call
    | let_assign
    | if conditional
    | member_access_expr
    | expr := expr
    | expr = expr
    | expr @ expr
    | expr @@ expr

func_asign -> function ID ( args_list ) => expr
    | function ID ( args_list ) : ID => expr
    | function ID ( args_list ) { lines }
    | function ID ( args_list ) : ID { lines }

let_assign -> let var_assign_list in expr
    | let var_assign_list in lines_block

// ... (resto de la gramática)
\end{verbatim}

\section{Ejemplos de Código}

\subsection{Hola Mundo}
\begin{verbatim}
print("Hello, World!");
\end{verbatim}

\subsection{Declaración de Variables y Expresiones}
\begin{verbatim}
let message = "El resultado es: " in {
    let a = 10;
    let b = 32;
    print(message @@ (a + b));
};
\end{verbatim}

\subsection{Definición de Funciones}
\begin{verbatim}
function add(a: Number, b: Number): Number => a + b;

print(add(5, 7));
\end{verbatim}

\subsection{Estructuras Condicionales}
\begin{verbatim}
let a = 10 in
    if (a > 5) {
        print("a es mayor que 5");
    } else {
        print("a no es mayor que 5");
    }
\end{verbatim}

\end{document}
